\documentclass[10pt,a4paper]{amsart}
\usepackage[margin=19mm]{geometry}
\usepackage{amsmath,amssymb,mathtools}
\usepackage[scr=rsfs,cal=euler]{mathalfa}
\usepackage{xfrac}
\usepackage[unicode,hidelinks,bookmarks=false]{hyperref}
\newcommand{\ie}{i.e.\ }
\newcommand{\cf}{cf.\ }
\newcommand{\resp}{respectively\ }
\newcommand{\Subsection}{\subsection}
\providecommand{\nobreakdash}{\nobreak\hbox{-}}
\setlength{\emergencystretch}{2em}
\multlinegap0pt
\usepackage{mathtools}
\usepackage[all]{xy}
\usepackage{enumerate}
\usepackage{amssymb}
\usepackage{tikz}
\newtheorem{lemma}{Lemma}
\DeclareMathOperator{\id}{id}
\title{$K$-regularity for finite quotient singularities}
\author{Matthew Satriano, Wanchun Shen}
\date{Source: arXiv:2609.21043v1 (CC BY 4.0)}
\begin{document}
\maketitle

\noindent\emph{Source and licence.} $K$-regularity for finite quotient singularities. Version arXiv:2609.21043v1 (CC BY 4.0), distributed by arXiv under the Creative Commons Attribution 4.0 International licence, http://creativecommons.org/licenses/by/4.0/, as recorded in the arXiv licence metadata for this exact version and shown on the abstract page https://arxiv.org/abs/2609.21043v1. Full licence notice: \texttt{licenses/arxiv-2609.21043-CC-BY-4.0.txt}.\par\smallskip

\noindent\textit{Editorial scope.} Counted unit: the article's lemma lem:induction (source0.tex lines 648--651). The article's complete original proof of it is delivered with it (source0.tex lines 652--670, one \texttt{proof} environment of 19 lines). No mathematics is written, repaired or re-derived: the statement and the proof are the article's own lines, in the article's own notation, and no retained line is substituted or re-flowed.  No further source-local context is needed: every cross-reference of every retained line resolves inside the two delivered spans.  Nothing was omitted from any retained span, so the package carries no omission record.  No retained line contains a citation, so the wrapper carries no bibliography and no bibliography entry.  No retained line contains a figure environment, an image-inclusion command or drawing code, so no image is delivered and the package declares no asset dependency.  Editorial formatting only, in addition: an AMS \texttt{amsart} preamble that restates the article's own declarations and macros used by the retained lines, namely the environments \texttt{lemma} on separate counters, together with the standard packages those lines need.  The retained environments print in this document as Lemma 1 in order of appearance, whereas the article prints the same items with the numbering of its own full document.  The complete native source of the article as distributed in the arXiv e-print for this exact version is bundled unchanged as \texttt{sources/210930/source0.tex}.
\begin{lemma}[Induction]\label{lem:induction}
Let $H\leq G$, let $W$ be an $H$-representation, and let
$V=\operatorname{Ind}_H^G W$.  If $C_H(W)\neq0$, then $C_G(V)\neq0$.
\end{lemma}

\begin{proof}
    The inclusion $W\subset V$ is $H$-equivariant, so we have an induced map $W/H\to V/G$ making the diagram
    \[
        \xymatrix{
            W\ar[d]_-{\pi_W}\ar[r] & V\ar[d]^-{\pi_V}\\
            W/H\ar[r]^-{i} & V/G
        }
    \]
    commute. We then obtain a commutative diagram where the rows are exact
    \[
        \xymatrix{
            i^*\Omega^1_{V/G}\ar[r]\ar[d] & i^*((\pi_V)_*\Omega^1_V)^G\ar[d]^-{\alpha}\ar[r] & i^*C_G(V)\ar[r]\ar[d] & 0\\
            \Omega^1_{W/H}\ar[r] & ((\pi_W)_*\Omega^1_W)^H\ar[r] & C_H(W)\ar[r] & 0
        }
    \]
    Upon showing $\alpha$ is surjective, we are done; indeed, we would then have $i^*C_G(V)\to C_H(W)$ is surjective so $i^*C_G(V)\neq0$ and hence $C_G(V)\neq0$.

    To prove $\alpha$ is surjective, it suffices to show $(\Omega^1_{k[V]})^G\to(\Omega^1_{k[W]})^H$ is surjective. Let $V=\bigoplus_{gH\in G/H}gW$, $p_{gH}\colon V\to gW$ be the projection, and $t_{gH}\colon W\xrightarrow{\simeq} gW$ be the isomorphism given by $t_{gH}(w)=gw$. If $\eta\in(\Omega^1_{k[W]})^H$, let $\eta_{gH}:=p_{gH}^*(t_{gH}^{-1})^*\eta$ and let $\eta':=\sum_{gH\in G/H}\eta_{gH}$. Then $\eta'\in (\Omega^1_{k[V]})^G$ and we see $i^*\eta'=\eta$. Indeed, $i^*p_{H}^*=\id$ and $i^*p_{gH}^*=0$ for $gH\neq H$. This proves surjectivity of $(\Omega^1_{k[V]})^G\to(\Omega^1_{k[W]})^H$, hence also of $\alpha$.
\end{proof}

\end{document}
